\documentclass[12pt]{article}
\usepackage[utf8]{inputenc}
\usepackage[russian]{babel}
\usepackage{amsmath}
\usepackage{amssymb}
\usepackage{geometry}
\geometry{a4paper, margin=2cm}

\begin{document}

\section*{Задача 3.1.7}

\textbf{Условие:} Алгоритм быстрого возведения в степень.\
\texttt{in:} $n$\\
\texttt{out:} $n^n$

\section*{Псевдокод}
Источник: https://github.com/cp-algorithms/cp-algorithms/blob/main/src/algebra/binary-exp.md
\begin{verbatim}
long long binpow(long long base, long long exp) {
    long long res = 1;
    while (exp > 0) {
        if (exp % 2 == 1)
            res *= base;
        base *= base;
        exp /= 2;
    }
    return res;
}
\end{verbatim}

\section*{Решение}

\begin{verbatim}
; Data Model
;0: r0
;1: result
;2: base
;3: exp

Load =0
Read 2
Sub 2
JZero END   ; 0^0 - invalid input
JGTZ END    ; (-a)^{-a} = 1 / (-a)^a = 0|1|-1. Suppoze n < 0 - invalid input

Load =1
Store 1     ; result := 1
Load 2
Store 3     ; exp := n

WHILE: Load =0
    Sub 3
    JZero END
    JGTZ END    ; if exp <= 0 (-exp >= 0) then return result
    Load 3      ; r0 := exp
    Div =2
    Mult =2
    Sub 3
    Mult =-1    ; r0 := ((exp // 2) * 2 - exp) * (-1) = exp % 2
    Sub =1
    JZero MULT_RES  ; if exp%2==1 then result *= base
    Jump WHILE_END
MULT_RES: Load 1
    Mult 2
    Store 1     ; result *= base
WHILE_END: Load 2
    Mult 2
    Store 2     ; base *= base
    Load 3
    Div =2
    Store 3     ; exp //= 2
    Jump WHILE
END: Write 1
    Halt 0
\end{verbatim}

\section*{Сложность}

Допустимый вход задачи - $n\geq 1$. Для $n\leq 0$ программа заканчивает работу без ответа.

Пусть
\[
 k:=\lfloor\log_2 n\rfloor+1
\]
- длина двоичной записи $n$, а $p$ - количество единиц в его двоичной записи.

Цикл WHILE проводится пока $exp>0$ и на каждом шаге exp //= 2, то есть shr exp. Значит цикл будет пройден $k$ раз. Заход в тело IF будет происходить $p$ раз. Значит худшим случаем будет $p=k$, то есть когда $n_{10}=1\dots1_2$.

\begin{enumerate}
    \item Проверка входа - $5$ шагов
    \item Инициализация - $4$ шага
    \item Тело цикла выполняется $k$ раз
    \begin{enumerate}
        \item На каждой итерации проверка условия - $4$ шага, (exp \% 2 == 1) - $7$ шагов
        \item Так как $p=k$, то (exp \% 2 == 1) всегда. Значит + 3 шага от MULT\_RES
        \item WHILE\_END - $7$ шагов
    \end{enumerate}
    \item Финальная проверка условия цикла - 3 шага
    \item END - 2 шага
\end{enumerate}

Значит, сложность в худшем случае:
\[
    T_{\max}(n) = 5+4+k(4 + 7 + 3 + 7) + 3 + 2 = 21k + 14 = 21(\lfloor\log_2 n\rfloor+1) + 14
\]

\end{document}